
\section{The pi-channel: discrete monodromy and cyclic closure}
\label{sec:pi_constraint}

In continuous settings, ``monodromy'' and ``single-valuedness'' are expressed via loops and holonomy.
Here we use the word \emph{monodromy} only as a mnemonic for a discrete \emph{closure} test:
a finite word determines a directed path on a constraint graph, and closure is detected by a wrap-around admissibility condition.
For clarity, the mathematical object of interest in this paper is \emph{cyclic closure} (periodic-orbit compatibility), not a continuous holonomy.

\subsection{The golden mean graph}
The golden mean shift admits a two-state Markov presentation.
Let the vertex set be $V:=\{0,1\}$ and define the allowed directed edges
\[
E:=\{0\to 0,\ 0\to 1,\ 1\to 0\},
\]
excluding $1\to 1$.
The adjacency matrix is
\[
A=\begin{pmatrix}1&1\\[2pt] 1&0\end{pmatrix}.
\]
Every word $w=w_1\cdots w_6\in X_6$ determines a directed path
\[
w_1\to w_2\to\cdots\to w_6
\]
of length $5$ on this graph.

\subsection{Endpoint defect via a boundary operator}
Let $\ZZ V$ be the free abelian group generated by $V$, and $\ZZ E$ the free abelian group generated by $E$.
Define the boundary operator $\partial:\ZZ E\to \ZZ V$ by
\[
\partial(u\to v):=v-u.
\]
For $w\in X_6$, let $c(w)\in \ZZ E$ be the path chain obtained by summing the edges traversed by the path $w_1\to\cdots\to w_6$.
Then
\[
\partial c(w)=w_6-w_1\in \ZZ V.
\]
We call $w$ \emph{endpoint-closed} if $\partial c(w)=0$, i.e.\ $w_1=w_6$, and \emph{endpoint-open} otherwise.

\subsection{Cyclic admissibility and the \texorpdfstring{$18\oplus 3$}{18+3} split}
For symbolic dynamics and zeta-function purposes, the relevant notion of closure is \emph{cyclic}:
a length-$6$ word corresponds to a periodic orbit segment only if the wrap-around transition $w_6\to w_1$ is also allowed.
For the golden mean constraint, this wrap-around admissibility is equivalent to excluding $w_6w_1=11$.

\begin{remark}[Why cyclic closure is the $\pi$-channel]
The cyclic closure condition is the one compatible with periodic-orbit counting and the Artin--Mazur zeta identities used in Section~\ref{sec:e_constraint}.
Endpoint closure ($w_1=w_6$) is also natural from a chain-boundary viewpoint, but it does not encode periodic admissibility for the golden mean constraint.
\end{remark}

\begin{definition}[Cyclic admissibility and boundary states]
Define
\[
X_6^{\mathrm{cyc}}:=\{w\in X_6:\ w_6w_1\neq 11\},
\qquad
X_6^{\mathrm{bdry}}:=X_6\setminus X_6^{\mathrm{cyc}}.
\]
We call elements of $X_6^{\mathrm{cyc}}$ \emph{cyclically admissible}, and elements of $X_6^{\mathrm{bdry}}$ \emph{boundary states}.
\end{definition}

\begin{proposition}[The boundary set has size $3$]
\label{prop:boundary_3}
One has
\[
X_6^{\mathrm{bdry}}=\{100001,\ 100101,\ 101001\},
\qquad |X_6^{\mathrm{bdry}}|=3,
\qquad |X_6^{\mathrm{cyc}}|=18.
\]
\end{proposition}

\begin{proof}
The wrap-around transition is forbidden if and only if $w_1=w_6=1$.
If $w_1=w_6=1$ and $w\in X_6$, then $w_2=w_5=0$.
The middle block $(w_3,w_4)$ can be $(0,0)$, $(0,1)$, or $(1,0)$ (but not $(1,1)$).
This yields exactly three boundary words:
\[
100001,\quad 100101,\quad 101001.
\]
Hence $|X_6^{\mathrm{bdry}}|=3$ and $|X_6^{\mathrm{cyc}}|=|X_6|-3=21-3=18$.
\end{proof}

\subsection{General length-n counts and the Lucas trace}
The cyclic/boundary split is not special to length $6$; it admits closed formulas at every length.

\begin{proposition}[Cyclic words, boundary words, and a Lucas-number identity]
\label{prop:cyclic_boundary_general}
Let $X_n\subset\{0,1\}^n$ be the set of length-$n$ words with no adjacent ones, and define
\[
X_n^{\mathrm{cyc}}:=\{w\in X_n:\ w_n w_1\neq 11\},
\qquad
X_n^{\mathrm{bdry}}:=X_n\setminus X_n^{\mathrm{cyc}}.
\]
Then for $n\ge 2$,
\[
|X_n^{\mathrm{bdry}}|=F_{n-2},
\qquad
|X_n^{\mathrm{cyc}}|=|X_n|-|X_n^{\mathrm{bdry}}|=F_{n+2}-F_{n-2}.
\]
Moreover, if $A=\bigl(\begin{smallmatrix}1&1\\ 1&0\end{smallmatrix}\bigr)$, then
\[
|X_n^{\mathrm{cyc}}|=\Tr(A^n)=F_{n-1}+F_{n+1},
\]
which equals the Lucas number $L_n$ (defined by $L_0=2$, $L_1=1$, and $L_{n+2}=L_{n+1}+L_n$).
\end{proposition}

\begin{proof}
A word $w\in X_n$ lies in $X_n^{\mathrm{bdry}}$ if and only if $w_1=w_n=1$.
Then admissibility forces $w_2=w_{n-1}=0$, and the middle block $(w_3,\dots,w_{n-2})$ is an arbitrary admissible word of length $n-4$.
Thus $|X_n^{\mathrm{bdry}}|=|X_{n-4}|=F_{(n-4)+2}=F_{n-2}$.
The identity $|X_n^{\mathrm{cyc}}|=|X_n|-|X_n^{\mathrm{bdry}}|$ is immediate.

For the trace formula, cyclic admissible words of length $n$ are in bijection with closed walks of length $n$ on the Markov graph with adjacency matrix $A$, hence their number is $\Tr(A^n)$; see, e.g., \cite{LindMarcus1995}.
To compute $\Tr(A^n)$, extend Fibonacci by $F_0:=0$ and note the matrix identity
\[
A^n=\begin{pmatrix}F_{n+1} & F_n\\[2pt] F_n & F_{n-1}\end{pmatrix}
\qquad (n\ge 1),
\]
which follows by induction using $A^{n+1}=A^n A$ and the Fibonacci recursion.
Taking the trace yields $\Tr(A^n)=F_{n+1}+F_{n-1}$.
\end{proof}

\begin{remark}[Two natural $\pi$-defects]
The $\pi$-channel admits two closely related nonnegative defects on $X_6$:
\begin{itemize}
\item \textbf{Endpoint defect:} $\mathcal{D}_\pi^{\mathrm{end}}(w):=\ind_{\{w_1\neq w_6\}}$.
\item \textbf{Cyclic defect:} $\mathcal{D}_\pi^{\mathrm{cyc}}(w):=\ind_{\{w_1=w_6=1\}}$.
\end{itemize}
The cyclic defect is the one compatible with periodic-orbit counting and zeta identities (Section~\ref{sec:e_constraint}), and it is responsible for the canonical $18\oplus 3$ split within $X_6$.
\end{remark}


